% Transcription of Cayley's Collected Mathematical Papers, Volume 4
% PDF pages 201--225 (printed pages 179--203)
% Papers: 258, 259, 260

\documentclass[11pt]{article}

% ============================================================
% PREAMBLE
% ============================================================
\usepackage[utf8]{inputenc}
\usepackage[T1]{fontenc}
\usepackage[english]{babel}
\usepackage{amsmath,amssymb}
\usepackage{geometry}
\usepackage{fancyhdr}
\usepackage{hyperref}
\usepackage[protrusion=true,expansion=false]{microtype}

\geometry{a4paper,margin=1in}

\pagestyle{fancy}
\fancyhf{}
\fancyhead[LE,RO]{\thepage}
\renewcommand{\headrulewidth}{0pt}

% Fallback providecommands
\providecommand{\mathbbm}[1]{\mathbb{#1}}
\providecommand{\sgn}{\operatorname{sgn}}

\setlength{\abovedisplayskip}{4pt plus 1pt minus 1pt}
\setlength{\belowdisplayskip}{4pt plus 1pt minus 1pt}

% Cayley delimiter convention shorthand
\newcommand{\Cdel}{\mid}

\hypersetup{
  pdftitle={Cayley Collected Papers Vol IV, pp.~179--203},
  pdfauthor={Arthur Cayley},
  hidelinks
}

\begin{document}

\thispagestyle{empty}
\begin{center}
\vspace*{2cm}
{\Large\scshape The Collected Mathematical Papers}\\[0.4em]
{\large of}\\[0.4em]
{\Large\scshape Arthur Cayley}\\[1.5em]
{\large Volume IV --- pages 179--203}\\[0.6em]
{\normalsize Papers 258, 259, 260}\\[2em]
\vfill
\end{center}
\newpage

\setcounter{page}{179}

% ============================================================
% PAPER 258
% ON A RELATION BETWEEN TWO TERNARY CUBIC FORMS
% ============================================================

\noindent\textbf{258]}\hfill\textbf{179}

\bigskip

\begin{center}
{\large\bfseries 258.}

\medskip

{\scshape On a Relation between two Ternary Cubic Forms.}
\end{center}

\medskip

\noindent[\textit{From the Philosophical Magazine}, vol.~xx.~1860, pp.~512--514.]

\bigskip

\noindent The cubic form
\[
x^3 + y^3 + z^3 + 6lxyz
\]
is in general linearly transformable into the form
\[
(X + Y + Z)^3 + 27kXYZ;
\]
in fact, writing
\[
\begin{aligned}
X &= 2lx - y - z,\\
Y &= 2ly - z - x,\\
Z &= 2lz - x - y,
\end{aligned}
\]
we have identically
\[
(1 - 2l + 4l^2)(X + Y + Z)^3 + 24(1 - l)^3 XYZ = 8(2l + 1)^3 (1 - l)^3 (x^3 + y^3 + z^3 + 6lxyz);
\]
and the value of $k$ consequently is
\[
k = -\frac{8(l-1)^3}{9(1 - 2l + 4l^2)}.
\]
If, however, $l = \tfrac{1}{2}$ or $l = -\tfrac{1}{2}$, the transformation fails. In the former case, viz.\ for $l = \tfrac{1}{2}$, the equations for the linear transformation become
\[
\begin{aligned}
X &= 2x - y - z,\\
Y &= 2y - z - x,\\
Z &= 2z - x - y,
\end{aligned}
\]
\hfill 23--2

\newpage

\noindent\textbf{180}\hfill\textsc{on a relation between two ternary cubic forms.}\hfill\textbf{[258}

\medskip

\noindent which give $X + Y + Z = 0$, so that $X$, $Y$, $Z$ are no longer independent; and the formula of transformation becomes
\[
(X + Y + Z)^3 = 0.
\]

It may be noticed that the invariant $S$ of the form
\[
x^3 + y^3 + z^3 + 6lxyz
\]
is $S = -l + l^4$, so that $l = 1$ is one of the values which make $S$ vanish. And the above transformation is not applicable to the cubic form $x^3 + y^3 + z^3 + 6xyz$, which is a form for which $S$ vanishes. The transformation, however, holds good for $l = 0$, which is another value which makes $S$ vanish; or it does apply to the form $x^3 + y^3 + z^3$, for which $S$ vanishes. The transformation, in fact, is
\[
(X + Y + Z)^3 + 24 XYZ = -8(x^3 + y^3 + z^3),
\]
with the linear equations
\[
\begin{aligned}
X &= -y - z,\\
Y &= -z - x,\\
Z &= -x - y.
\end{aligned}
\]

The above two forms for which $S$ vanishes, viz.
\[
x^3 + y^3 + z^3 + 6xyz,\qquad x^3 + y^3 + z^3,
\]
are, notwithstanding, equivalent to each other, as appears by the identical equation
\[
(x + y + z)^3 + (x + \omega y + \omega^2 z)^3 + (x + \omega^2 y + \omega z)^3 = 3(x^3 + y^3 + z^3 + 6xyz),
\]
where $\omega$ is an imaginary cube root of unity. In the latter of the two cases of failure, viz.\ for $l = -\tfrac{1}{2}$, the equations for the linear transformations are
\[
X = Y = Z = -x - y - z;
\]
so that $X$, $Y$, $Z$ are not only not independent, but they are connected by two linear relations. And the formula of transformation becomes
\[
(X + Y + Z)^3 - 27 XYZ = 0,
\]
which is, in fact, true in virtue of the equations $X = Y = Z$.

The two forms of equation,
\[
x^3 + y^3 + z^3 + 6lxyz = 0,\qquad (x + y + z)^3 + 27k\,xyz = 0,
\]
represent each of them equally well a curve of the third order without a double point. In the first form the three real points of inflexion are given by
\[
(x = 0,\ y + z = 0),\quad (y = 0,\ z + x = 0),\quad (z = 0,\ x + y = 0);
\]

\newpage

\noindent\textbf{258]}\hfill\textsc{on a relation between two ternary cubic forms.}\hfill\textbf{181}

\medskip

\noindent or what is the same thing, the points in question are the intersections of the lines $x = 0$, $y = 0$, $z = 0$ with the line $x + y + z = 0$; or we have $x + y + z = 0$ for the equation of the line through the three points of inflexion; and the equations of the tangents at the points of inflexion are
\[
2lx - y - z = 0,\quad 2ly - z - x = 0,\quad 2lz - x - y = 0.
\]
For the second form it is obvious that the points of inflexion are the intersections of the lines $x = 0$, $y = 0$, $z = 0$ with the line $x + y + z = 0$; and, moreover, that the lines $x = 0$, $y = 0$, $z = 0$ are the tangents at the point of inflexion.

The first of the above-mentioned forms, however, cannot represent a curve with a double point. In fact the condition for its doing so would be $1 + 8l^3 = 0$; but when this condition is satisfied, the left-hand side breaks up into linear factors, and the equation represents, not a proper curve of the third order, but a system of three lines. The second form can represent a curve having a double point; viz.\ if $k = -1$, the curve will have a conjugate or isolated point at the point $x = y = z$. It is clear \textit{a priori} that ($x = 0$, $y = 0$, $z = 0$ being real lines) neither of the forms can represent a curve of the third order having a double point with two real branches through it, since in this case the curve has only one real point of inflexion.

I have elsewhere used the word ``node'' to denote a double point, and I take the opportunity of suggesting the employment of the words ``crunode'' (\textit{crus}) and ``acnode'' (\textit{acus}) to denote respectively a double point with two real branches through it, and a conjugate or isolated point.

\medskip

\noindent 2, \textit{Stone Buildings, W.C., October} 19, 1860.

\newpage

% ============================================================
% PAPER 259
% THE PROBLEM OF POLYHEDRA
% ============================================================

\noindent\textbf{182}\hfill\textbf{[259}

\bigskip

\begin{center}
{\large\bfseries 259.}

\medskip

{\scshape The Problem of Polyhedra.}
\end{center}

\medskip

\noindent[\textit{From the Philosophical Transactions of the Royal Society of London}, vol.~\textsc{cxlvii}.~(for the year 1857), pp.~183--185: printed as a Note to Mr Kirkman's Memoir ``On Autopolar Polyhedra,'' pp.~183--215.]

\bigskip

\noindent Let $a$, $b$, $c$, $d$, $e$, $f$, $g$, $h$, \&c.\ represent the vertices of a polyhedron, then a face will be represented, e.g.\ by $abcde$, where the contiguous duads, viz.\ $ab$, $bc$, $cd$, $de$, $ea$ are the edges of the face; and calling the face $K$, we may write
\[
K = abcde. \tag{1}
\]

It is to be noticed that the letters of a face-symbol may be taken forwards or backwards from any letter without altering the meaning of the symbol. Thus, $abcde$, $bcdea$, \&c., $edcba$, \&c.\ might any of them be taken to denote the face $K$. The diagonal of a face cannot be either an edge or a diagonal of any other face, i.e.\ a non-contiguous duad such as $ac$ in a face-symbol $K$ cannot be a duad, contiguous or non-contiguous, of any other face-symbol. But each edge of a face must be an edge of one and only one other face, i.e.\ each contiguous duad such as $ab$ in the face-symbol $K$ must be a contiguous duad of one and only one other face-symbol $L$. And moreover two faces cannot have more than a single edge in common, i.e.\ two face-symbols cannot contain more than a single contiguous duad, the same in each symbol.

The face $K$ contains the edges $ab$, $ae$, i.e.\ the edge $ab$ is contained in the face $K$; it will also be contained in one and only one other face, suppose $L$; this face will contain another edge through the vertex $a$, suppose the edge $af$, and so on, until we arrive at a face containing the edge $ae$; we have, for example,

\newpage

\noindent\textbf{259]}\hfill\textsc{the problem of polyhedra.}\hfill\textbf{183}

\medskip

\[
\begin{aligned}
K &= eabcd,\\
L &= baf\ldots,\\
M &= fag\ldots,\\
N &= gah\ldots,\\
P &= hai\ldots,\\
Q &= iae\ldots;
\end{aligned}
\]
and we thence derive the vertex-symbol
\[
a = KLMNPQ, \tag{2}
\]
where the contiguous duads $KL$, $LM$, $MN$, $NP$, $PQ$, $QK$ represent in order the edges through the vertex $a$. The remarks before made with respect to the face-symbols apply to the vertex-symbols. A non-contiguous duad such as $KN$ of the vertex-symbol $a$ cannot be a duad, contiguous or non-contiguous, of any other vertex-symbol; but each contiguous duad such as $KL$ of the vertex-symbol $a$ must be a contiguous duad of one and only one other vertex-symbol $b$. And the symbols of two vertices cannot contain more than one contiguous duad, the same in each symbol.

Any edge of the polyhedron admits of a double representation; it is the junction of two vertices, or the intersection of two faces. Thus $ab$ and $KL$ will represent the same edge, or we may write
\[
ab = KL. \tag{3}
\]

It is to be remarked that in this system, to each equation $ab = KL$ there corresponds one and only one equation of the form $ae = KQ$, i.e.\ to an edge considered as drawn from a given vertex in a given face there corresponds one and only one other edge from the same vertex in the same face.

It has been shown how the system of face-symbols (1) leads to the system of vertex-symbols (2), and the system of edge symbols (3); and generally, any one of the three systems leads to the other two; and the three systems conjointly, or each system by itself, is a complete representation of the polyhedron. As an example, take the hexaedron: the three systems are:
\[
\begin{array}{lll}
K = abcd, & a = LPK, & (2)\\
L = abfe, & b = LMK, & \\
M = bfgc, & c = MNK, & \\
N = gcdh, & d = NPK, & \\
P = dhea, & e = LPQ, & \\
Q = hefg, & f = LMQ, & \\
& g = MNQ, & \\
& h = NPQ; &
\end{array}
\]

\newpage

\noindent\textbf{184}\hfill\textsc{the problem of polyhedra.}\hfill\textbf{[259}

\medskip

\[
\begin{array}{lll}
ab = KL, & ae = PL, & ef = LQ, \quad (3)\\
bc = KM, & bf = LM, & fg = MQ, \\
cd = KN, & cg = MN, & gh = NQ, \\
de = KP, & dh = NP, & he = PQ.
\end{array}
\]

\medskip

\begin{center}
[\textit{Figure: schematic of the hexaedron showing the labelling of faces $K$, $L$, $M$, $N$, $P$, $Q$ and vertices $a$, $b$, $c$, $d$, $e$, $f$, $g$, $h$.}]
\end{center}

\medskip

Consider, now, two polyhedra having the same number of vertices and also the same number of faces. And let the vertices and faces of the first polyhedron taken in any order be represented by
\[
abcde\ldots KLM\ldots,
\]
and the vertices and faces of the second polyhedron taken in a certain order be represented by
\[
a'b'c'd'e'\ldots K'L'M'\ldots;
\]
then, forming the substitution symbol
\[
a'b'c'd'e'\ldots K'L'M'\ldots abcde\ldots KLM\ldots,
\]
which denotes that $a'$ is to be written for $a$, $b'$ for $b$\ldots $K'$ for $K$, \&c., if operating with this upon the symbol system of the first polyhedron, we obtain the symbol system of the second polyhedron, the second polyhedron will be syntypic with the first. It should be noticed, that there may be several modes of arrangement of the vertices and faces of the second polyhedron, which will render it syntypic according to the foregoing definition with the first polyhedron, i.e.\ the second polyhedron may be syntypic in several different ways with the first polyhedron. This is, in fact, the same as saying that a polyhedron may be syntypic with itself in several different ways. Suppose, next, that the number of vertices of the second polyhedron is equal to the number of faces of the first polyhedron, and the number of faces of the second polyhedron is equal to the number of vertices of the first polyhedron: and let the vertices and faces of the first polyhedron in any order be represented by
\[
abcde\ldots KLM\ldots,
\]
and the faces and vertices of the second polyhedron in a certain order be represented by
\[
A'B'C'D'E'\ldots k'l'm'\ldots.
\]

\newpage

\noindent\textbf{259]}\hfill\textsc{the problem of polyhedra.}\hfill\textbf{185}

\medskip

\noindent Then, forming the substitution symbol
\[
A'B'C'D'E'\ldots k'l'm'\ldots abcde\ldots KLM\ldots,
\]
if, operating with this upon the symbol system of the first polyhedron, we obtain the symbol system of the second polyhedron, the second polyhedron is said to be polar-syntypic with the first; and, as in the case of syntypicism, this may happen in several different ways.

Lastly, if there be a polyhedron having the same number of vertices and faces, and if the vertices and faces in any order be represented by
\[
abcd\ldots KLMN\ldots,
\]
and the faces and vertices in a certain order be represented by
\[
ABCD\ldots klmn\ldots;
\]
then, forming the substitution symbol
\[
ABCD\ldots klmn\ldots abcd\ldots KLMN\ldots,
\]
if, operating with this upon the symbol system of the polyhedron, we reproduce such symbol system, i.e.\ in fact, if the polyhedron be polar-syntypic with itself, the polyhedron is said to be autopolar; and in accordance with a preceding remark, this may happen in several different ways. It is clear that the substitution symbol, operating on the symbol system of the vertices, must give the symbol system of the faces, and conversely; but operating on the symbol system of the edges, it must reproduce such symbol system of the edges: and this last condition will by itself suffice to make the polyhedron autopolar, i.e.\ the polyhedron will be autopolar if the substitution symbol, operating on the symbol system of the edges, reproduces such symbol system.

\bigskip

\noindent\hfill C. IV. \hfill 24

\newpage

% ============================================================
% PAPER 260
% ON THE DOUBLE TANGENTS OF A PLANE CURVE
% ============================================================

\noindent\textbf{186}\hfill\textbf{[260}

\bigskip

\begin{center}
{\large\bfseries 260.}

\medskip

{\scshape On the Double Tangents of a Plane Curve.}
\end{center}

\medskip

\noindent[\textit{From the Philosophical Transactions of the Royal Society of London}, vol.~\textsc{cxlix}, for the year 1859, pp.~193--212. Received March 17,---Read April 14, 1859.]

\bigskip

\noindent It was first shown by Pl\"ucker on geometrical principles, that the number of the double tangents of a plane curve of the order $m$ was $\tfrac{1}{2}m(m-2)(m^2-9)$: see the note, ``Solution d'une question fondamentale concernant la th\'eorie g\'en\'erale des Courbes,'' \textit{Crelle}, t.~\textsc{xii}.~pp.~105--108 (1834), and the ``Theorie der algebraischen Curven'' (1839). The memoir by Hesse, ``Ueber die Wendepuncte der Curven dritter Ordnung,'' \textit{Crelle}, t.~\textsc{xxviii}.~pp.~97--107 (1844), contains the analytical solution of the allied easier problem of the determination of the points of inflexion of a plane curve. In the memoir, ``Recherches sur l'\'elimination et sur la th\'eorie des Courbes,'' \textit{Crelle}, t.~\textsc{xxxiv}.~(1847), pp.~30--45, [53], I showed how the problem of double tangents admitted of an analytical solution, viz.\ if $U = 0$ is the equation of the curve, $L$, $M$, $N$ the first derived functions of $U$, and
\[
D = \alpha(M\partial_z - N\partial_y) + \beta(N\partial_y - L\partial_z) + \gamma(L\partial_y - M\partial_x)
\]
(where $\alpha$, $\beta$, $\gamma$ are arbitrary), then the points of contact of the double tangents are given as the intersections of the curve $U = 0$, with a curve the equation whereof is in the first instance obtained under the form $[F] = 0$; $[F]$ being a given function of
\[
DU,\ D^2U,\ \ldots\ D^{m-1}U,
\]
of the degree $m^2 - m - 6$ in respect of $(\alpha, \beta, \gamma)$, the degree $m^3 - 2m^2 - 10m + 12$ in respect of $(x, y, z)$, and the degree $m^2 + m - 12$ in respect of the coefficients of $U$. It was necessary, in order that the points of intersection should be independent of the arbitrary quantities $(\alpha, \beta, \gamma)$, that we should have identically
\[
[F] = \Lambda \cdot U + N \cdot \Pi U,
\]

\newpage

\noindent\textbf{260]}\hfill\textsc{on the double tangents of a plane curve.}\hfill\textbf{187}

\medskip

\noindent $N$ being of the degree $m^2 - m - 6$ in $(\alpha, \beta, \gamma)$, and consequently $\Pi U$ a function of $(x, y, z)$ without $(\alpha, \beta, \gamma)$. Guided by Hesse's investigation for the points of inflexion, I asserted that it was probable that $N$ was of the form $(\alpha x + \beta y + \gamma z)^{m^2-m-6}$; which being so, $\Pi U$ would be of the degree $(m-2)(m^2-9)$ in respect of $(x, y, z)$, and the degree $m^2 + m - 12$ in respect of the coefficients, and I was thus led to the theorem, ``On trouve les points de contact des tangentes doubles en combinant avec l'\'equation de la courbe une \'equation $\Pi U = 0$, de l'ordre $(m-2)(m^2-9)$ par rapport aux variables et de l'ordre $m^2 + m - 12$ par rapport aux coefficients---c'est-\`a-dire, puisqu'il correspond deux points de contact \`a une tangente double, le nombre de ces tangentes est \'egal \`a $\tfrac{1}{2}m(m-2)(m^2-9)$: th\'eor\`eme d\'emontr\'e indirectement par M.~Pl\"ucker.''

Hesse, in the memoir ``Ueber Curven dritter Ordnung u.s.w.,'' \textit{Crelle}, t.~\textsc{xxxvi}.~pp.~143--176 (1848), showed how the components $D^iU$, $D^2U$, $\ldots$ $D^{m-1}U$ of $[F]$ could each of them be expressed in a simplified form, and he thus effected the actual reduction of $[F]$ to the form $\Lambda \cdot U + (\alpha x + \beta y + \gamma z)^{m-2} \cdot R$, where $R$ still contained the arbitrary quantities $(\alpha, \beta, \gamma)$ in the degree $(m-2)(m-3)$. In particular for a quartic curve, the equation $R = 0$ was shown to be
\[
(\mathfrak{A}, \mathfrak{B}, \mathfrak{C}, \mathfrak{F}, \mathfrak{G}, \mathfrak{H}\,\Cdel\alpha, \beta, \gamma)^2 \cdot H = 0,
\]
where the left-hand side is of the degree 2 in $(\alpha, \beta, \gamma)$ and the degree 16 in $(x, y, z)$; and which should therefore by means of the equation $U = 0$ be reducible so as to contain the factor $(\alpha x + \beta y + \gamma z)^2$.

Jacobi's paper, ``Beweis des Satzes, dass eine Curve $n$-ten Grades im allgemeinen $\tfrac{1}{2}n(n-2)(n^2-9)$ Doppeltangenten hat,'' \textit{Crelle}, t.~\textsc{xl}.~pp.~237--260 (1850), did not, I think, materially advance the solution of the question. In a letter to Jacobi, dated the 30th December, 1849, published at the conclusion of the last-mentioned paper, Hesse gave the equation of the curve of the 14th order for the points of contact of the double tangents of a quartic, viz.\ in my notation,
\[
(\mathfrak{A}, \mathfrak{B}, \mathfrak{C}, \mathfrak{F}, \mathfrak{G}, \mathfrak{H}\,\Cdel\partial_x H,\ \partial_y H,\ \partial_z H)^2 \cdot H - H\,(\mathfrak{A}, \mathfrak{B}, \mathfrak{C}, \mathfrak{F}, \mathfrak{G}, \mathfrak{H}\,\Cdel\partial_x,\ \partial_y,\ \partial_z)^2 H = 0,
\]
and the demonstration is given in Hesse's paper, ``Ueber die ganzen homogenen Functionen von der dritten und vierten Ordnung zwischen drei Variabeln,'' \textit{Crelle}, t.~\textsc{xli}.~pp.~285--292 (1851), and is reproduced in Mr Salmon's Treatise on the Higher Plane Curves (1852). Two very interesting memoirs by Hesse and Steiner, \textit{Crelle}, t.~\textsc{xlix}.~(1855), relate to the geometrical theory of the double tangents of a quartic, and it is not necessary to refer to them more particularly. It is to be observed that the curve which determines the points of contact of the double tangents is not absolutely determinate; for we may, it is clear, in the place of $\Pi U = 0$, write $\Pi U + M \cdot U = 0$, where $M$ is an arbitrary function of the proper degree: a very elegant transformation in the case of the quartic is given in Hesse's paper, ``Transformation der Gleichung der Curven 14ten Grades, welche eine gegebene Curve 4ten Grades in den Ber\"uhrungspuncten ihrer Doppeltangenten schneiden,'' \textit{Crelle}, t.~\textsc{lii}.~pp.~97--103 (1856).

Mr Salmon's work above referred to, contains the fundamental theorem of the tangential of a cubic, viz.\ a tangent to a cubic meets the cubic in a third point
\hfill 24--2

\newpage

\noindent\textbf{188}\hfill\textsc{on the double tangents of a plane curve.}\hfill\textbf{[260}

\medskip

\noindent which lies on the second or line polar of the point of contact with respect to the Hessian. In my ``Memoir on Curves of the Third Order,'' \textit{Phil.~Trans.}~vol.~\textsc{cxlvii}.~(1857), pp.~415--446, art.~No.~37, [146], I gave an identical equation relating to the tangential of a cubic, but which is not there exhibited in its proper form; this was afterwards effected by Mr Salmon, in the paper ``On Curves of the Third Order,'' \textit{Phil.~Trans.}~vol.~\textsc{cxlviii}.~(1858), pp.~535--541. The equation, as given by Mr Salmon, is in the notation of the present memoir,
\[
-\mathfrak{H} \cdot U + \tfrac{1}{3}\mathfrak{D}\mathfrak{H} \cdot DU - \tfrac{1}{3}DH \cdot \mathfrak{D}\mathbf{T} + H \cdot \mathbf{T} = 0,
\]
an equation which in fact puts in evidence the last-mentioned theorem for the tangential of a cubic.

The idea occurred to me of considering, in the case of the higher plane curves, the tangentials of a given point of the curve, viz.\ the points in which the tangent again meets the curve; for by expressing that two of these tangentials were coincident, we should have the condition that the given point is the point of contact of a double tangent. But I was not able to complete the solution.

Finally, Mr Salmon discovered the equation of a curve of the order $n - 2$, which by its intersections with the tangent at the given point determines the tangentials, and by expressing that the curve in question is touched by the tangent, he was led to a complete solution of the Double-tangent problem. Mr Salmon's result is given in the note, ``On the Double Tangents to Plane Curves,'' in the Philosophical Magazine for October 1858. The discovery just referred to led me to the investigations of the present memoir, in which it will be seen that I obtain, for a curve of any order whatever, the identical equation corresponding to the before-mentioned equation obtained by Mr Salmon in the case of a cubic; which identical equation puts in evidence the theorem as to the tangentials of the curve, and may thus be considered as containing in itself the solution of the Double-tangent problem: the identical equation is besides interesting for its own sake, as a part of the theory of ternary quantics.

\medskip

1. Mr Salmon's solution of the problem of double tangents is based upon the following analytical determination of the tangentials of any point of the curve.

Let
\[
T = (*\,\Cdel X, Y, Z)^n = 0
\]
be the equation of the given curve, $(X, Y, Z)$ being current coordinates; and let $(x, y, z)$ be the coordinates of a point on the curve, so that we have
\[
U = (*\,\Cdel x, y, z)^n = 0,
\]
a condition satisfied by the coordinates of the point in question.

Then the tangent
\[
V = (X\partial_x + Y\partial_y + Z\partial_z)U = 0
\]
at the point $(x, y, z)$, meets the curve besides in $(n-2)$ points, which are the tangentials of the given point $(x, y, z)$, and which are determined as the intersections of the tangent $V = 0$ with a certain curve,
\[
\Omega = (*\,\Cdel x, y, z)^{n-2} = 0.
\]

\newpage

\noindent\textbf{260]}\hfill\textsc{on the double tangents of a plane curve.}\hfill\textbf{189}

\medskip

\noindent 2. To express the equation of this curve, let $U_1$, $U_2$, $\ldots$ be the successive emanants of $U$, taken with the facients of emanation $(x_1, y_1, z_1)$, viz.
\[
U_1 = \frac{1}{n}\bigg[\text{should be } \frac{1}{n-2}\bigg](x_1\partial_x + y_1\partial_y + z_1\partial_z)U,
\]
\[
U_2 = \frac{1}{n(n-1)}\bigg[\text{should be } \frac{1}{(n-2)(n-3)}\bigg](x_1\partial_x + y_1\partial_y + z_1\partial_z)^2 U,
\]
\[
\ldots
\]
where it should be noticed that the numerical determination is such, that putting $(x, y, z)$ for $(x_1, y_1, z_1)$, then $U_1$, $U_2$, $\ldots$ become respectively equal to $U$. [The numerical determination should have been and in the latter part of the memoir is assumed to be such as to render $H_1$, $H_2$, \&c.\ equal to $H$, on making the substitution in question: the correction was made in a later memoir ``On the double tangents of a curve of the fourth order.''] Suppose also that $H$, $H_1$, $H_2$, $\ldots$ are the Hessians of $U$, $U_1$, $U_2$, $\ldots$, viz.\ $H$ is the determinant formed with the second derived functions of $U$ with respect to $(x, y, z)$, $H_1$ the like determinant with the second derived functions of $U_1$ with respect to the same quantities $(x, y, z)$; and so on. Moreover let $D^{n-2}H_r = (X\partial_x + Y\partial_y + Z\partial_z)^{n-2}H$, denote the $(n-2)$thic emanant of $H$ with respect to the current coordinates $(X, Y, Z)$ as facients of emanation; and similarly let $D^{n-2}H_1$, $D^{n-2}H_2$, $\ldots$ denote the $(n-2)$thic emanants of $H_1$, $H_2$, $\ldots$ in respect to the same facient of emanation---it being understood that in all these functions, $(x_1, y_1, z_1)$ are after the differentiations to be replaced by $(x, y, z)$. It is to be observed that $U_r$ is of the degree $(n-r)$ in $(x, y, z)$, and consequently $H_r$ of the degree $3(n-2-r)$; hence $D^{n-2}H_r$ is of the degree $3(n-2-r)-(n-2)$, $= 2(n-2)-3r$, which implies that $r \not> \tfrac{1}{3}(n-2)$, for otherwise $D^{n-2}H_r$ would be identically equal to zero. Upon replacing $(x_1, y_1, z_1)$ by $(x, y, z)$, $D^{n-2}H_r$ ($r$ satisfying the above condition) becomes of the degree $2(n-2)$ in $(x, y, z)$, and it is obviously of the degree 3 in the coefficients of $U$, and of the degree $(n-2)$ in the current coordinates $(X, Y, Z)$.

\medskip

3. This being premised, we have
\[
\Omega = D^{n-2}H - \frac{n-1}{1}D^{n-2}H_1 + \&c.\ = 0,
\]
for the equation of the curve of the order $(n-2)$, which by its intersection with the tangent gives the tangentials of the given point; the numerical coefficients are the binomial coefficients of the order $(n-1)$ taken with the signs $+$ and $-$ alternately, and the series is continued as long as the terms do not vanish, that is, if as before $r$ denote the suffix of $H$, for so long as $r \not> \tfrac{1}{3}(n-2)$; but of course the value will not be altered by continuing the series to $r = n - 1$. In particular, for the quartic we have
\[
\Omega = D^2H - 3D^2H_1,
\]
for the quintic
\[
\Omega = D^3H - 4D^3H_1 + 6D^3H_2,
\]

\newpage

\noindent\textbf{190}\hfill\textsc{on the double tangents of a plane curve.}\hfill\textbf{[260}

\medskip

\noindent and so on. The function $\Omega$, like the several component terms, is of course of the degree 3 in the coefficients of $U$, and of the degree $2(n-2)$ in $(x, y, z)$.

\medskip

4. It is to be remarked that the formula applies to a cubic; we have here simply $\Omega = DH$, which agrees with a result already mentioned. It may be noticed also that in the general case the formula gives at once the condition for the points of inflexion; in fact, if the point $(x, y, z)$ be a point of inflexion, then one of the tangentials must coincide with this point, or the equation $\Omega = 0$ will be satisfied by writing therein $(x, y, z)$ for $(X, Y, Z)$; but when this is done $D^{n-2}H$, $D^{n-2}H_1$, \&c.\ reduce themselves (to numerical factors \textit{pr\`es}) to $H$, and the equation becomes simply $H = 0$, which is the well-known condition for the points of inflexion.

\medskip

5. If two of the tangentials coincide, or what is the same thing, if the tangent $V = 0$ touches the curve $\Omega = 0$, then the point $(x, y, z)$ will be the point of contact of a double tangent. The equation which expresses the condition in question, treating therein $(x, y, z)$ as current coordinates, is consequently that of a curve, intersecting the given curve (now represented by $U = 0$) in the points of contact of the double tangents. The process leads to a determinate form $\Pi U = 0$, of the curve in question, but of course any curve whatever, $\Pi U + M \cdot U = 0$, will intersect the curve $U = 0$ in the points of contact of the double tangents.

\medskip

6. I write for the moment
\[
\Omega = (A, \ldots\Cdel X, Y, Z)^{n-2} = 0,
\]
\[
V = \xi X + \eta Y + \zeta Z = 0,
\]
for the two equations: the coefficients $(A, \ldots)$, as already mentioned, are of the degree $2(n-2)$ in $(x, y, z)$ and of the degree 3 in the coefficients of $U$; or as we may express it,
\[
A, \ldots = (a, \ldots)^3\,(x, y, z)^{2(n-2)}.
\]
In like manner $\xi$, $\eta$, $\zeta$ are of the degree $(n-1)$ in $(x, y, z)$, and the degree 1 in the coefficients of $U$, or we may write
\[
\xi, \eta, \zeta = (a, \ldots)^1\,(x, y, z)^{n-1}.
\]

\medskip

7. The equation which expresses that the line $V = 0$ touches the curve $\Omega = 0$, is $\nabla\Omega = 0$, where the facients of the Reciprocant $\nabla\Omega$ are the coefficients $(\xi, \eta, \zeta)$ of the linear function. This equation is of the form
\[
(A, \ldots)^{2(n-3)}\,(\xi, \eta, \zeta)^{(n-2)(n-3)} = 0;
\]
or attending to the forms of $(A, \ldots)$ and $(\xi, \eta, \zeta)$, it is of the form
\[
(a, \ldots)^{n+4(n-3)}\,(x, y, z)^{(n-2)\cdot 2(n-3) + (n-1)(n-2)(n-3)} = 0,
\]
or what is the same thing, the form
\[
(a, \ldots)^{n+4(n-3)}\,(x, y, z)^{(n-2)(n^2-9)} = 0,
\]

\newpage

\noindent\textbf{260]}\hfill\textsc{on the double tangents of a plane curve.}\hfill\textbf{191}

\medskip

\noindent viz.\ the curve through the points of contact of the double tangents is a curve of the order $(n-2)(n^2-9)$, and its equation contains the coefficients of the equation $U = 0$ of the given curve in the degree $(n+4)(n-3)$. And since each double tangent corresponds to two points of contact, the number of double tangents is $\tfrac{1}{2}n(n-2)(n^2-9)$. This agrees with the before-mentioned results.

\medskip

8. The whole problem is thus reduced to the demonstration of Mr Salmon's expression for the curve $\Omega = 0$. To fix the ideas, consider the case of a quartic curve $T = (*\,\Cdel X, Y, Z)^4 = 0$, and let the function $U = (*\,\Cdel x, y, z)^4$ (or as for shortness we may write it, $U = (x, y, z)^4$) and certain of its emanants be represented as follows, viz.---
\[
\begin{aligned}
a &= (x, y, z)^4,\\
b &= (x, y, z)^3(X, Y, Z),\\
c &= (x, y, z)^2(X, Y, Z)^2,\\
d &= (x, y, z)(X, Y, Z)^3,\\
e &= (X, Y, Z)^4,\\[4pt]
a' &= (x, y, z)^3(X', Y', Z'),\\
b' &= (x, y, z)^2(X, Y, Z)(X', Y', Z'),\\
c' &= (x, y, z)(X, Y, Z)^2(X', Y', Z'),\\
d' &= (X, Y, Z)^3(X', Y', Z'),\\[4pt]
a'' &= (x, y, z)^2(X', Y', Z')^2,\\
b'' &= (x, y, z)(X, Y, Z)(X', Y', Z')^2,\\
c'' &= (X, Y, Z)^2(X', Y', Z')^2,
\end{aligned}
\]
where $(X', Y', Z')$ are new arbitrary facients; but, as before, $(X, Y, Z)$ are taken to be current coordinates, and $(x, y, z)$ the coordinates of the given point on the curve:

$e = 0$ is the equation of the curve;

$d = 0$, the equation of the first or cubic polar of the point $(x, y, z)$;

$b = 0$, the equation of the last or line polar of the point $(x, y, z)$, or what is the same thing (the point being on the curve), the tangent of the curve at this point;

$a = 0$, the condition which expresses that the point is on the curve.

\medskip

9. Imagine now an identical equation,
\[
a\,\mathrm{I} + b\,\mathrm{II} + d\,\mathrm{III} + e\,\mathrm{IV} = 0;
\]
then, since $a = 0$, we have
\[
b\,\mathrm{II} + d\,\mathrm{III} + e\,\mathrm{IV} = 0;
\]

\newpage

\noindent\textbf{192}\hfill\textsc{on the double tangents of a plane curve.}\hfill\textbf{[260}

\medskip

\noindent and if in this equation we write $b = 0$, $e = 0$, it becomes $\mathrm{III}\,d = 0$, that is, the points of intersection of the curve $e = 0$ and the tangent $b = 0$ lie on one or other of the curves $d = 0$, $\mathrm{III} = 0$. But the points in question do not lie on the curve $d = 0$, consequently they lie on the curve $\mathrm{III} = 0$.

\medskip

10. To explain the law of formation of the multipliers $\mathrm{I}$, $\mathrm{II}$, $\mathrm{III}$, $\mathrm{IV}$, I form the matrix
\[
\begin{pmatrix}
a, & b, & c, & d, & e; & a', & b', & c', & d';\\
b, & c, & d, & e, & f; & b', & c', & d', & e';\\
a', & b', & c', & d', & e'; & a'', & b'', & c'', & d''
\end{pmatrix}
\]
\textit{(matrix layout reconstructed from scan; the relevant rows for the quartic are the top two and the accented row, with the appropriate truncation.)}

\medskip

and then we have
\[
\mathrm{I} = -\begin{vmatrix} d, & e, & b'\\ e, & f, & c'\\ d', & e', & b'' \end{vmatrix},\quad
\mathrm{II} = +\begin{vmatrix} d, & e, & a'\\ e, & f, & b'\\ d', & e', & a'' \end{vmatrix} - \begin{vmatrix} c, & e, & b'\\ d, & f, & c'\\ c', & e', & b'' \end{vmatrix},
\]
\[
\mathrm{III} = -\begin{vmatrix} b, & c, & b'\\ c, & d, & c'\\ b', & c', & b'' \end{vmatrix} - \begin{vmatrix} a, & c, & a'\\ b, & d, & b'\\ a', & c', & a'' \end{vmatrix},\quad
\mathrm{IV} = +\begin{vmatrix} a, & b, & a'\\ b, & c, & b'\\ a', & b', & a'' \end{vmatrix},
\]
values which, as I proceed to show, satisfy the identical equation
\[
a\,\mathrm{I} + b\,\mathrm{II} + d\,\mathrm{III} + e\,\mathrm{IV} = 0.
\]

\medskip

11. We have in fact
\[
\begin{aligned}
\mathrm{I} = {}& d(db'' - c'^2 + cc'' - b'd')\\
& + e(b'c' - b''c + b'c' - bc'')\\
& + d'(cc' - db' + bd' - cc'),
\end{aligned}
\]
where the last line is $= bd'^2 - db'd'$;
\[
\begin{aligned}
\mathrm{II} = {}& d(b'c' - a''d + b'c' - b''c + a'd' - bc'')\\
& + e(ca'' - a'c' + bb'' - b'^2 + ac'' - a'c')\\
& + d'(a'd - b'c + b'c - bc' + bc' - ad'),
\end{aligned}
\]
where the last line is $da'd' - ad'^2$;

\newpage

\noindent\textbf{260]}\hfill\textsc{on the double tangents of a plane curve.}\hfill\textbf{193}

\medskip

\[
\begin{aligned}
\mathrm{III} = {}& a(-b'd' + cc'' + c'^2 - b''d + b'd' - ea'')\\
& + b(bc'' - b'c' + b''c - b'c' + a''d - ad')\\
& + a'(cc' - bd' + b'd - cc' + ae - b'd),
\end{aligned}
\]
where the last line is $-ba'd' + ea'^2$; and
\[
\begin{aligned}
\mathrm{IV} = {}& a(b''c - b'c' + da'' - b'c')\\
& + b(b'^2 - bb'' + a'c' - a''c)\\
& + a'(bc' - b'c + cb' - a'd),
\end{aligned}
\]
where the last line is $ba'c - a'^2 d$.

These values may be expressed as follows:
\[
\begin{array}{rl}
\mathrm{I} = & a(\ 0\ )\\
& {} + b(\ d'^2 - \ \ \ )\quad (12)\\
& {} + d(\ b''d + c''c - b'd' - cc' - d'b'\ )\quad (13)\\
& {} + e(-b''c - c''b + b'c' + c'b'\ )\quad (14)\\[4pt]
\mathrm{II} = & a(-d'^2\ )\quad (21)\\
& {} + b(\ 0\ )\\
& {} + d(-a''d - b''c - c''b + a'd' + b'c' + c'b' + d'a')\quad (23)\\
& {} + e(\ a''c + b''b + c''a - a'c' - b'b' - c'a'\ )\quad (24)\\[4pt]
\mathrm{III} = & a(-b''d - c''c + b'd' + c'c' + d'b'\ )\quad (31)\\
& {} + b(\ a''d + b''c + c''b - a'd' - b'c' - c'b' - d'a')\quad (32)\\
& {} + d(\ 0\ )\\
& {} + e(-a''a + a'a'\ )\quad (34)\\[4pt]
\mathrm{IV} = & a(\ b''c + c''b - b'c' - c'b'\ )\quad (41)\\
& {} + b(-a''c - b''b - c''a + a'c' + b'b' + c'a'\ )\quad (42)\\
& {} + d(\ a''a - a'a'\ )\quad (43)\\
& {} + e(\ 0\ ).
\end{array}
\]

\noindent\hfill C. IV. \hfill 25

\newpage

\noindent\textbf{194}\hfill\textsc{on the double tangents of a plane curve.}\hfill\textbf{[260}

\medskip

\noindent which are of the form
\[
\begin{aligned}
\mathrm{I} &= a\cdot 0 + b(12) + d(13) + e(14),\\
\mathrm{II} &= a(21) + b\cdot 0 + d(23) + e(24),\\
\mathrm{III} &= a(31) + b(32) + d\cdot 0 + e(34),\\
\mathrm{IV} &= a(41) + b(42) + d(43) + e\cdot 0,
\end{aligned}
\]
where $(12) = -(21)$, \&c., and which therefore satisfy the equation
\[
a\,\mathrm{I} + b\,\mathrm{II} + d\,\mathrm{III} + e\,\mathrm{IV} = 0.
\]

\medskip

12. The equation of the curve which by its intersection with the tangent gives the tangentials, is
\[
\mathrm{III} = -\begin{vmatrix} b, & c, & b'\\ c, & d, & c'\\ b', & c', & b'' \end{vmatrix} - \begin{vmatrix} a, & c, & a'\\ b, & d, & b'\\ a', & c', & a'' \end{vmatrix},
\]
the degrees of which are
\[
\begin{aligned}
&\text{in the coefficients of } U,\quad 3,\\
&\text{in } (x, y, z),\quad 6,\\
&\text{in } (X, Y, Z),\quad 4,\\
&\text{in } (X', Y', Z'),\quad 2;
\end{aligned}
\]
and it only remains to divest this equation of a factor which it contains,
\[
\begin{vmatrix} x, & y, & z\\ X, & Y, & Z\\ X', & Y', & Z' \end{vmatrix}^2,
\]
which being thrown out, the equation will be independent of $(X', Y', Z')$ and will be of the degrees
\[
\begin{aligned}
&\text{in the coefficients of } U,\quad 3,\\
&\text{in } (x, y, z),\quad 4,\\
&\text{in } (X, Y, Z),\quad 2.
\end{aligned}
\]
and will in fact be the before-mentioned equation $\Omega = D^2H - 3D^2H_1 = 0$.

\medskip

13. Write for shortness,
\[
\Delta = \begin{vmatrix} x, & y, & z\\ X, & Y, & Z\\ X', & Y', & Z' \end{vmatrix};
\]
it is to be shown that
\[
\mathrm{III} = -\tfrac{1}{6}\Delta^2(D^2H - 3D^2H_1).
\]

\newpage

\noindent\textbf{260]}\hfill\textsc{on the double tangents of a plane curve.}\hfill\textbf{195}

\medskip

\noindent 14. To effect this I remark that we have identically
\[
\begin{vmatrix} a, & b, & a'\\ b, & c, & b'\\ a', & b', & a'' \end{vmatrix} = \Delta^2 H;
\]
and I proceed to operate upon this equation with $D = X\partial_x + Y\partial_y + Z\partial_z$.

I notice that
\[
a,\ b,\ c,\ d,\ e;\quad a',\ b',\ c',\ d';\quad a'',\ b'',\ c''
\]
are in regard to $(x, y, z)$ of the degrees
\[
4,\ 3,\ 2,\ 1,\ 0;\quad 3,\ 2,\ 1,\ 0;\quad 2,\ 1,\ 0;
\]
or what is the same thing, since for the case in hand $n = 4$, of the degrees
\[
n,\ n-1,\ \ldots;\quad n-1,\ n-2,\ \ldots;\quad n-2,\ n-3,\ \ldots
\]
and we have
\[
Da = nb,\quad Db = (n-1)c,\ldots\quad Da' = (n-1)b',\quad Db' = (n-2)c',\ldots\quad Da'' = (n-2)b'',\ldots
\]

\medskip

15. In the determinant
\[
\begin{vmatrix} a, & b, & a'\\ b, & c, & b'\\ a', & b', & a'' \end{vmatrix} = \Delta^2 H,
\]
the degrees of the terms (other than each top term, the degree of which is higher by unity) in the several columns are $n-1$, $n-2$, $n-2$; if then we operate on the determinant with $D$, and as regards the top terms we write
\[
\begin{aligned}
Da &= b + (n-1)b,\\
Db &= c + (n-2)c,\\
Da' &= b' + (n-2)b',
\end{aligned}
\]
we have in the first place a term
\[
\begin{vmatrix} b, & b, & b'\\ c, & c, & c'\\ b', & b', & a'' \end{vmatrix} + \begin{vmatrix} a, & c, & a'\\ b, & d, & b'\\ a', & c', & a'' \end{vmatrix} + \begin{vmatrix} a, & b, & b'\\ b, & c, & c'\\ a', & b', & b'' \end{vmatrix}
\]
next the terms
\[
(n-1)\begin{vmatrix} b, & b, & a'\\ c, & c, & b'\\ b', & b', & a'' \end{vmatrix} + (n-2)\begin{vmatrix} a, & c, & a'\\ b, & d, & b'\\ a', & c', & a'' \end{vmatrix} + (n-2)\begin{vmatrix} a, & b, & b'\\ b, & c, & c'\\ a', & b', & b'' \end{vmatrix},
\]
\hfill 25--2

\newpage

\noindent\textbf{196}\hfill\textsc{on the double tangents of a plane curve.}\hfill\textbf{[260}

\medskip

\noindent the first of which vanishes. On the right-hand side $D\Delta = 0$ identically, and therefore $D\cdot\Delta^2 H = \Delta^2 DH$, or we have
\[
(n-2)\begin{vmatrix} a, & c, & a'\\ b, & d, & b'\\ a', & c', & a'' \end{vmatrix} + (n-2)\begin{vmatrix} a, & b, & b'\\ b, & c, & c'\\ a', & b', & b'' \end{vmatrix} = \Delta^2 DH.
\]

\medskip

16. I repeat the operation $D$: we have
\[
\begin{aligned}
&(n-2)(n-1)\begin{vmatrix} a, & d, & a'\\ b, & e, & b'\\ a', & d', & a'' \end{vmatrix} + (n-2)(n-1)\begin{vmatrix} b, & c, & a'\\ c, & d, & b'\\ b', & c', & a'' \end{vmatrix}\\
&{} + (n-2)(n-3)\begin{vmatrix} a, & c, & b'\\ b, & d, & c'\\ a', & c', & b'' \end{vmatrix} + (n-2)(n-2)\begin{vmatrix} a, & b, & c'\\ b, & c, & d'\\ a', & b', & c'' \end{vmatrix}\\
&{} + (n-2)(n-2)\begin{vmatrix} b, & b, & b'\\ c, & c, & c'\\ b', & b', & b'' \end{vmatrix} + (n-2)(n-3)\begin{vmatrix} a, & b, & b'\\ b, & c, & c'\\ a', & b', & b'' \end{vmatrix} = \Delta^2 D^2 H,
\end{aligned}
\]
or, collecting the different terms,
\[
(n-2)(n-3)\begin{vmatrix} a, & c, & b'\\ b, & d, & c'\\ a', & c', & b'' \end{vmatrix} + 2(n-2)^2\begin{vmatrix} a, & b, & c'\\ b, & c, & d'\\ a', & b', & c'' \end{vmatrix} + (n-2)(n-3)\begin{vmatrix} a, & b, & b'\\ b, & c, & c'\\ a', & b', & b'' \end{vmatrix} + (n-2)(n-1)\begin{vmatrix} a, & d, & a'\\ b, & e, & b'\\ a', & d', & a'' \end{vmatrix} = \Delta^2 D^2 H.
\]

\medskip

17. A little consideration will show that in this equation we may write $n - 1$ for $n$, and $H_1$ for $H$. In fact, putting for a moment $\delta = x_1\partial_x + y_1\partial_y + z_1\partial_z$, corresponding to the equation
\[
\begin{vmatrix} a, & b, & a'\\ b, & c, & b'\\ a', & b', & a'' \end{vmatrix} = \Delta^2 H,
\]
this other equation,
\[
\begin{vmatrix} \delta a, & \delta b, & \delta a'\\ \delta b, & \delta c, & \delta b'\\ \delta a', & \delta b', & \delta a'' \end{vmatrix} = \Delta^2 H_1,
\]
where ultimately $(x_1, y_1, z_1)$ are to be replaced by $(x, y, z)$. We may operate upon this equation with $D$, $D^2$, $\ldots$ as before, the only difference being that in the first

\newpage

\noindent\textbf{260]}\hfill\textsc{on the double tangents of a plane curve.}\hfill\textbf{197}

\medskip

\noindent instance $\delta a$, $\delta b$, \&c.\ are as regards $(x, y, z)$ of degrees lower by unity than $a$, $b$, \&c., that is $n - 1$ must be substituted throughout in the place of $n$; and when at the end of the process $(x_1, y_1, z_1)$ are replaced by $(x, y, z)$, then $\delta a$, $\delta b$, \&c.\ become equal to $a$, $b$, \&c., from which the truth of the asserted proposition is manifest.

\medskip

18. Hence writing $n = 4$, we have
\[
2\begin{vmatrix} a, & b, & c'\\ b, & c, & d'\\ a', & b', & c'' \end{vmatrix} + 8\begin{vmatrix} a, & d, & a'\\ b, & e, & b'\\ a', & d', & a'' \end{vmatrix} + 2\begin{vmatrix} a, & c, & b'\\ b, & d, & c'\\ a', & c', & b'' \end{vmatrix} + 6\begin{vmatrix} b, & c, & a'\\ c, & d, & b'\\ b', & c', & a'' \end{vmatrix} = \Delta^2 D^2 H,
\]
\[
2\begin{vmatrix} a, & d, & a'\\ b, & e, & b'\\ a', & d', & a'' \end{vmatrix} + 2\begin{vmatrix} b, & c, & a'\\ c, & d, & b'\\ b', & c', & a'' \end{vmatrix} = \Delta^2 D^2 H_1,
\]
and hence
\[
2\begin{vmatrix} a, & b, & c'\\ b, & c, & d'\\ a', & b', & c'' \end{vmatrix} + 2\begin{vmatrix} a, & c, & b'\\ b, & d, & c'\\ a', & c', & b'' \end{vmatrix} + 2\begin{vmatrix} b, & c, & a'\\ c, & d, & b'\\ b', & c', & a'' \end{vmatrix} = \Delta^2(D^2H - 3D^2H_1),
\]
which is the required equation
\[
\mathrm{III} = -\tfrac{1}{6}\Delta^2(D^2H - 3D^2H_1).
\]

\medskip

19. It is to be added, that the equation for $\Delta^2 DH$ gives $\mathrm{IV} = \tfrac{1}{2}\Delta^2 DH$; the values of $\mathrm{II}$ and $\mathrm{I}$ are at once obtained from those of $\mathrm{III}$ and $\mathrm{IV}$ by interchanging $(x, y, z)$ and $(X, Y, Z)$. Hence if we represent by $\mathfrak{H}$, $\mathfrak{D}$, \&c.\ the values which $H$, $D$, \&c.\ assume by this interchange, we may write
\[
\begin{aligned}
\mathrm{I} &= -\tfrac{1}{2}\Delta^2\mathfrak{D}\mathfrak{H},\\
\mathrm{II} &= +\tfrac{1}{6}\Delta^2(\mathfrak{D}^2\mathfrak{H} - 3\mathfrak{D}^2\mathfrak{H}_1),\\
\mathrm{III} &= -\tfrac{1}{6}\Delta^2(D^2H - 3D^2H_1),\\
\mathrm{IV} &= +\tfrac{1}{2}\Delta^2 DH;
\end{aligned}
\]
and the identical equation,
\[
a\,\mathrm{I} + b\,\mathrm{II} + d\,\mathrm{III} + e\,\mathrm{IV} = 0,
\]
gives therefore
\[
-\mathfrak{D}\mathfrak{H} \cdot U + \tfrac{1}{3}(\mathfrak{D}^2\mathfrak{H} - 3\mathfrak{D}^2\mathfrak{H}_1)\,DU - \tfrac{1}{3}(D^2H - 3D^2H_1)\,\mathfrak{D}\mathbf{T} + DH \cdot \mathbf{T} = 0,
\]
which is of itself sufficient to put in evidence the property that the curve $D^2H - 3D^2H_1 = 0$ gives by its intersections with the tangent $DU = 0$, the tangentials of the point $(x, y, z)$. The last-mentioned equation is the equation for a quartic corresponding to Mr Salmon's equation
\[
-\mathfrak{H} \cdot U + \tfrac{1}{3}\mathfrak{D}\mathfrak{H} \cdot DU - \tfrac{1}{3}DH \cdot \mathfrak{D}\mathbf{T} + H \cdot \mathbf{T} = 0
\]
for the cubic $U = 0$.

\newpage

\noindent\textbf{198}\hfill\textsc{on the double tangents of a plane curve.}\hfill\textbf{[260}

\medskip

\noindent 20. It is worth while to give the investigation of the equation for the cubic; the matrix is
\[
\begin{pmatrix}
a, & b, & c, & d;\ \  a', & b'\\
b, & c, & d, & e;\ \  b', & c'\\
a', & b', & c', & d';\ \ a'', & b''
\end{pmatrix}
\]
and the identical equation is
\[
a\,\mathrm{I} + b\,\mathrm{II} + c\,\mathrm{III} + d\,\mathrm{IV} = 0,
\]
where
\[
\begin{aligned}
\mathrm{I} &= \begin{vmatrix} c, & b, & b'\\ d, & c, & c'\\ c', & b', & b'' \end{vmatrix} - \begin{vmatrix} a, & c, & a'\\ b, & d, & b'\\ a', & c', & a'' \end{vmatrix},\\[4pt]
\mathrm{II} &= -\begin{vmatrix} c, & b, & a'\\ d, & c, & b'\\ c', & b', & a'' \end{vmatrix},\\[4pt]
\mathrm{III} &= -\begin{vmatrix} a, & b, & a'\\ b, & c, & b'\\ a', & b', & a'' \end{vmatrix},\\[4pt]
\mathrm{IV} &= \begin{vmatrix} a, & b, & a'\\ b, & c, & b'\\ a', & b', & a'' \end{vmatrix},
\end{aligned}
\]
\textit{(matrix variants reconstructed from scan: the III and IV expressions follow the same pattern with the appropriate column substitutions.)}

or, as we may express them,
\[
\begin{array}{rl}
\mathrm{I} = & a(\ 0\ )\\
& {} + b(\ c'^2\ )\quad (12)\\
& {} + c(\ b''c - b'c' - c'b'\ )\quad (13)\\
& {} + d(-b''b + b'b'\ )\quad (14)\\[4pt]
\mathrm{II} = & a(-c'^2\ )\quad (21)\\
& {} + b(\ 0\ )\\
& {} + c(-a''c - b''b + a'c' + b'b' + c'a')\quad (23)\\
& {} + d(\ a''b + b''a - a'b' - b'a'\ )\quad (24)\\[4pt]
\mathrm{III} = & a(-b''c + b'c' + c'b'\ )\quad (31)\\
& {} + b(\ a''c + b''b - a'c' - b'b' - c'a')\quad (32)\\
& {} + c(\ 0\ )\\
& {} + d(-a''a + a'a'\ )\quad (34)\\[4pt]
\mathrm{IV} = & a(\ b''b - b'b'\ )\quad (41)\\
& {} + b(-a''b - b''a + a'b' + b'a'\ )\quad (42)\\
& {} + c(\ a''a - a'a'\ )\quad (43)\\
& {} + d(\ 0\ )
\end{array}
\]

\newpage

\noindent\textbf{260]}\hfill\textsc{on the double tangents of a plane curve.}\hfill\textbf{199}

\medskip

\noindent which verify the identical equation. We have $\mathrm{III} = -\Delta^2 DH$, $\mathrm{IV} = \Delta^2 H$, and thence $\mathrm{II} = +\Delta^2\mathfrak{D}\mathfrak{H}$, $\mathrm{I} = -\Delta^2\mathfrak{H}$; hence the equation in question
\[
-\mathfrak{H} \cdot U + \tfrac{1}{3}\mathfrak{D}\mathfrak{H} \cdot DU - \tfrac{1}{3}DH \cdot \mathfrak{D}\mathbf{T} + H \cdot \mathbf{T} = 0.
\]

\medskip

21. One other example will be sufficient to render manifest the law of the formation of the multipliers $\mathrm{I}$, $\mathrm{II}$, $\mathrm{III}$, $\mathrm{IV}$.

In the case of a sextic curve we have the matrix
\[
\begin{pmatrix}
a, & b, & c, & d, & e, & f;\ \  a', & b', & c', & d', & e'\\
b, & c, & d, & e, & f, & g;\ \  b', & c', & d', & e', & f'\\
a', & b', & c', & d', & e', & f';\ \ a'', & b'', & c'', & d'', & e''
\end{pmatrix}
\]
the identical equation is
\[
a\,\mathrm{I} + b\,\mathrm{II} + f\,\mathrm{III} + g\,\mathrm{IV} = 0;
\]
and the expressions for the multipliers $\mathrm{I}$, $\mathrm{II}$, $\mathrm{III}$, $\mathrm{IV}$ are:
\[
\begin{aligned}
\mathrm{I} = & \begin{vmatrix} f, & e, & b'\\ g, & f, & c'\\ f', & e', & b'' \end{vmatrix} + \begin{vmatrix} f, & d, & c'\\ g, & e, & d'\\ f', & d', & c'' \end{vmatrix} + \begin{vmatrix} f, & c, & d'\\ g, & d, & e'\\ f', & c', & d'' \end{vmatrix} + \begin{vmatrix} f, & b, & e'\\ g, & c, & f'\\ f', & b', & e'' \end{vmatrix},\\[6pt]
\mathrm{II} = & -\begin{vmatrix} f, & e, & a'\\ g, & f, & b'\\ f', & e', & a'' \end{vmatrix} - \begin{vmatrix} f, & d, & b'\\ g, & e, & c'\\ f', & d', & b'' \end{vmatrix} - \begin{vmatrix} f, & c, & c'\\ g, & d, & d'\\ f', & c', & c'' \end{vmatrix} - \begin{vmatrix} f, & b, & d'\\ g, & c, & e'\\ f', & b', & d'' \end{vmatrix} - \begin{vmatrix} f, & a, & e'\\ g, & b, & f'\\ f', & a', & e'' \end{vmatrix},\\[6pt]
\mathrm{III} = & -\begin{vmatrix} a, & b, & e'\\ b, & c, & f'\\ a', & b', & e'' \end{vmatrix} - \begin{vmatrix} a, & c, & d'\\ b, & d, & e'\\ a', & c', & d'' \end{vmatrix} - \begin{vmatrix} a, & d, & c'\\ b, & e, & d'\\ a', & d', & c'' \end{vmatrix} - \begin{vmatrix} a, & e, & b'\\ b, & f, & c'\\ a', & e', & b'' \end{vmatrix} - \begin{vmatrix} a, & f, & a'\\ b, & g, & b'\\ a', & f', & a'' \end{vmatrix},\\[6pt]
\mathrm{IV} = & \begin{vmatrix} a, & b, & d'\\ b, & c, & e'\\ a', & b', & d'' \end{vmatrix} + \begin{vmatrix} a, & c, & c'\\ b, & d, & d'\\ a', & c', & c'' \end{vmatrix} + \begin{vmatrix} a, & d, & b'\\ b, & e, & c'\\ a', & d', & b'' \end{vmatrix} + \begin{vmatrix} a, & e, & a'\\ b, & f, & b'\\ a', & e', & a'' \end{vmatrix}.
\end{aligned}
\]

\medskip

22. We have in fact
\[
\begin{array}{rl}
\mathrm{I} = & a(\ 0\ )\\
& {} + b(-ge'' + ff'\ )\quad (12)\\
& {} + f(\ b''f + c''e + d''d + e''c - b'f' - c'e' - d'd' - e'c' - f'b'\ )\quad (13)\\
& {} + g(-b''e - c''d - d''c - e''b + b'e' + c'd' + d'c' + e'b'\ )\quad (14)\\[4pt]
\mathrm{II} = & a(\ e''g\ )\quad (21)\\
& {} + b(\ 0\ )\\
& {} + f(-a''f - b''e - c''d - d''c - e''b + a'f' + b'e' + c'd' + d'c' + e'b' + f'a')\quad (23)\\
& {} + g(\ a''e + b''d + c''c + d''b - a'e' - b'd' - c'c' - d'b' - e'a'\ )\quad (24)
\end{array}
\]

\newpage

\noindent\textbf{200}\hfill\textsc{on the double tangents of a plane curve.}\hfill\textbf{[260}

\medskip

\[
\begin{array}{rl}
\mathrm{III} = & a(-b''f - c''e - d''d - e''c + b'f' + c'e' + d'd' + e'c' + f'b'\ )\quad (31)\\
& {} + b(\ a''f + b''e + c''d + d''c + e''b - a'f' - b'e' - c'd' - d'c' - e'b' - f'a')\quad (32)\\
& {} + f(\ 0\ )\\
& {} + g(-a''a + a'a'\ )\quad (34)\\[4pt]
\mathrm{IV} = & a(\ b''e + c''d + d''c + e''b - b'e' - c'd' - d'c' - e'b'\ )\quad (41)\\
& {} + b(-a''e - b''d - c''c - d''b + a'e' + b'd' + c'c' + d'b' + e'a')\quad (42)\\
& {} + f(\ a''a - a'a'\ )\quad (43)\\
& {} + g(\ 0\ )
\end{array}
\]
which are of the form
\[
\begin{aligned}
\mathrm{I} &= a\cdot 0 + b(12) + f(13) + g(14),\\
\mathrm{II} &= a(21) + b\cdot 0 + f(23) + g(24),\\
\mathrm{III} &= a(31) + b(32) + f\cdot 0 + g(34),\\
\mathrm{IV} &= a(41) + b(42) + f(43) + g\cdot 0,
\end{aligned}
\]
where $(12) = -(21)$ \&c., and the equation
\[
a\,\mathrm{I} + b\,\mathrm{II} + f\,\mathrm{III} + g\,\mathrm{IV} = 0
\]
is consequently satisfied.

\medskip

23. The expression
\[
\mathrm{III} = -\begin{vmatrix} a, & b, & e'\\ b, & c, & f'\\ a', & b', & e'' \end{vmatrix} - \begin{vmatrix} a, & c, & d'\\ b, & d, & e'\\ a', & c', & d'' \end{vmatrix} - \begin{vmatrix} a, & d, & c'\\ b, & e, & d'\\ a', & d', & c'' \end{vmatrix} - \begin{vmatrix} a, & e, & b'\\ b, & f, & c'\\ a', & e', & b'' \end{vmatrix} - \begin{vmatrix} a, & f, & a'\\ b, & g, & b'\\ a', & f', & a'' \end{vmatrix}
\]
leads to
\[
\mathrm{III} = -\Delta^2(D^4 H - 5D^4 H_1 + 10D^4 H_2),
\]
and consequently the equation of the curve which by its intersections with the tangent determines the tangentials of a point of a sextic, is
\[
D^4 H - 5D^4 H_1 + 10D^4 H_2 = 0.
\]

\medskip

24. In the general case of a curve of the order $n$ the matrix is
\[
\begin{pmatrix}
a_0, & a_1, & a_2, & \ldots & a_{n-1};\ \  a'_0, & a'_1, & \ldots & a'_{n-2}\\
a_1, & a_2, & a_3, & \ldots & a_n;\ \      a'_1, & a'_2, & \ldots & a'_{n-1}\\
a'_0, & a'_1, & a'_2, & \ldots & a'_{n-1};\ \ a''_0, & a''_1, & \ldots & a''_{n-2}
\end{pmatrix}
\]
where, in analogy with what precedes,
\[
\begin{aligned}
a_0 &= (X, Y, Z)^n,\\
a_1 &= (X, Y, Z)^{n-1}(x, y, z),\\
&\ \ \vdots\\
a_{n-1} &= (X, Y, Z)(x, y, z)^{n-1},\\
a_n &= (x, y, z)^n,
\end{aligned}
\]

\newpage

\noindent\textbf{260]}\hfill\textsc{on the double tangents of a plane curve.}\hfill\textbf{201}

\medskip

\noindent and similarly for the accented letters, so that

$a_0 = 0$ is the equation of the curve;

$a_1 = 0$ is the equation of the first or $(n-1)$thic polar;

$a_{n-1} = 0$ is the equation of the last or line-polar, or what is the same thing, since $(x, y, z)$ is a point on the curve, the tangent at this point;

$a_n = 0$, the condition which expresses that $(x, y, z)$ is a point of the curve;

and we have to form the identical equation
\[
a_0\,\mathrm{I} + a_1\,\mathrm{II} + a_{n-1}\,\mathrm{III} + a_n\,\mathrm{IV} = 0.
\]

\medskip

25. If, for shortness, the columns of the last-mentioned matrix are represented by
\[
1,\ 2,\ 3\ \ldots\ n,\ (1),\ (2)\ldots(n-1),
\]
and the determinants formed with these columns respectively by a corresponding notation $\{1, 2, (1)\}$, $\{1, 2, (2)\}$, \&c., then the expressions for the multipliers $\mathrm{I}$, $\mathrm{II}$, $\mathrm{III}$, $\mathrm{IV}$ are as follows, viz.
\[
\begin{aligned}
\mathrm{I} &= \{n,\ n-1,\ (2)\} + \{n,\ n-2,\ (3)\}\ \ldots\ + \{n,\ 2,\ (n-1)\},\\
\mathrm{II} &= -\{n,\ n-1,\ (1)\} - \{n,\ n-2,\ (2)\}\ \ldots\ - \{n,\ 2,\ (n-2)\} - \{n,\ 1,\ (n-1)\},\\
\mathrm{III} &= -\{1,\ 2,\ (n-1)\} - \{1,\ 3,\ (n-2)\}\ \ldots\ - \{1,\ n-1,\ (2)\} - \{1,\ n,\ (1)\},\\
\mathrm{IV} &= \{1,\ 2,\ (n-2)\} + \{1,\ 3,\ (n-1)\}\ \ldots\ + \{1,\ n-1,\ (1)\};
\end{aligned}
\]
the truth of the identical equation being shown, as in the foregoing special cases, by the transformation of the multipliers into the form
\[
\begin{aligned}
\mathrm{I} &= a_0 \cdot 0 + a_1(12) + a_{n-1}(13) + a_n(14),\\
\mathrm{II} &= a_0(21) + a_1 \cdot 0 + a_{n-1}(23) + a_n(24),\\
\mathrm{III} &= a_0(31) + a_1(32) + a_{n-1} \cdot 0 + a_n(34),\\
\mathrm{IV} &= a_0(41) + a_1(42) + a_{n-1}(43) + a_n \cdot 0,
\end{aligned}
\]
where $(12) = -(21)$, \&c.: the required expressions may be written down without difficulty.

\medskip

26. Proceeding then to reduce the equation
\[
\mathrm{III} = -\{1,\ 2,\ (n-1)\} - \{1,\ 3,\ (n-2)\}\ \ldots\ - \{1,\ n-1,\ (2)\} - \{1,\ n,\ (1)\},
\]
we have the equation
\[
\{1,\ 2,\ (1)\} = \Delta^2 H,
\]
which is to be successively operated on with $D$. The degrees (less unity) of the columns
\[
1,\ 2,\ \ldots\ n-1,\ n,\ (1)\ (2),\ \ldots\ n-1,
\]
are
\[
n-1,\ n-2,\ \ldots\ 1,\ 0,\ n-2,\ n-3,\ \ldots\ 0;
\]
and the rule is to operate on each column of the determinant, multiplying by the degree less unity, and increasing the symbolical number by unity. Thus
\[
D\{1, 2, (1)\} = (n-1)\{2, 2, (1)\} + (n-2)\{1, 3, (1)\} + (n-2)\{1, 2, (2)\},
\]
\[
= (n-2)\{1, 3, (1)\} + (n-2)\{1, 2, (2)\},
\]
since $\{2, 2, (1)\}$ vanishes identically. The following Table shows the mode of effecting the operations:

\noindent\hfill C. IV. \hfill 26

\newpage

\noindent\textbf{202}\hfill\textsc{on the double tangents of a plane curve.}\hfill\textbf{[260}

\medskip

\noindent\textit{[The original page contains an extensive table headed by $H$, $DH$, $D^2H$, $D^3H$, $D^4H$, $D^5H$, $D^6H$, $D^7H$, etc., showing the numerical coefficients (with factorials of $[n-2]$ and $[n-1]$) which multiply the symbolic determinants $\{1\,2\,(r)\}$, $\{1\,3\,(r)\}$, $\{1\,4\,(r)\}$, \&c.\ in the expansion of $\Delta^2 D^r H$. The general law of the table is summarised in art.~27 below.]}

\medskip

\noindent The first three columns of the table show the numbers which give, by the addition of the numbers in the same horizontal line, the numerical coefficients of the factorials which multiply the different terms of $H$, $DH$, \&c., and where in the last column $12(1)$, \&c.\ are written for shortness in the place of $\{1, 2, (1)\}$, \&c.

\newpage

\noindent\textbf{260]}\hfill\textsc{on the double tangents of a plane curve.}\hfill\textbf{203}

\medskip

\noindent where the first three columns show the numbers which give, by the addition of the numbers in the same horizontal line, the numerical coefficients of the factorials which multiply the different terms of $H$, $DH$, \&c., and where in the last column $12(1)$, \&c.\ are written for shortness in the place of $\{1, 2, (1)\}$, \&c.

\medskip

27. It is clear that we have in general
\[
\begin{aligned}
\Delta^2 D^r H = {}& \frac{1}{1}\,[n-2]^r[n-2]^r\,\{1, 2, (r+1)\}\\
&{} + \frac{r}{1}\,[n-2]^r[n-2]^{r-1}\,\{1, 3, (r)\}\\
&{} + \frac{r\cdot r - 1}{1\cdot 2}\,[n-2]^r[n-2]^{r-2}\,\{1, 4, (r-1)\}\\
&{} + \quad\vdots\\
&{} + \frac{1}{1}\,[n-2]^r[n-2]^0\,\{1, r+2, (1)\}\\[4pt]
&{} + [n-1]^1\,\left\{\frac{r}{1} + R'_{r-1}\right\}[n-2]^{r-1}[n-2]^{r-2}\,\{2, 3, (r-1)\}\\
&\qquad{} + \cdots\\
&\qquad{} + R'^{r-1}_{r-1}[n-2]^{r-1}[n-2]^0\,\{2, r+1, (1)\}\\[4pt]
&{} + [n-1]^2\,\left\{\frac{r}{1} + R''_{r-2}\right\}[n-2]^{r-2}[n-2]^{r-3}\,\{3, 4, (r-3)\}\\
&\qquad{} + \cdots\\
&\qquad{} + R''^{r-2}_{r-2}[n-2]^{r-2}[n-2]^0\,\{3, r, (1)\}\\[4pt]
&{} + \cdots\\[4pt]
&{} + [n-1]^{\lfloor (r-1)/2\rfloor}\,\left[\begin{array}{c}\text{$r$ even or odd}\end{array}\right]\,[n-2]^{\lceil r/2\rceil}\,[\cdots]\,\{\tfrac{1}{2}r+1,\ \tfrac{1}{2}r+2,\ (1)\}\\[4pt]
&{} + [n-1]^{\lfloor r/2\rfloor}\,[n-2]^{r/2}\,[n-2]^0\,\left\{\tfrac{1}{2}(r+1),\ \tfrac{1}{2}(r+3),\ (2)\right\}\\
&{} + [n-1]^{\lfloor r/2\rfloor + 1}\,[n-2]^{r/2 + 1}\,\left\{\tfrac{1}{2}(r+1),\ \tfrac{1}{2}(r+5),\ (1)\right\}.
\end{aligned}
\]
and the general term is
\[
[n-1]^\delta R_r^\delta\,[n-2]^{\beta+1}[n-2]^{r-2\delta - s}\,\{\beta+1,\ \beta+2+s,\ (r-2\delta+1-s)\},
\]
where $s$ extends from $s = 0$ to $s = r - 2\delta$, and $\delta$ from $\delta = 0$ to $\delta = \tfrac{1}{2}r$ or $\tfrac{1}{2}(r-1)$, according as $r$ is even or odd. The expression for the coefficients $R_r^\delta$ is
\[
R_r^\delta = \frac{[r]^\delta}{[\delta]^\delta},
\]
and that of the other coefficients $R_{r-s}^\delta$ ($\delta = $ or $< 1$) is not required for the present purpose.

\medskip

28. According to a remark already made, the expressions for $D^{n-2}H_1$, $D^{n-2}H_2$, \&c.\ are at once obtained from that for $D^rH$ by merely writing $n-1$, $n-2$, \&c.\ in the place of $n$: it is however to be noticed, that the quantity within the $[\ ]$ must not be negative, and that on its becoming so, the factorial is to be omitted.
\hfill 26--2

\end{document}
